A massa do tubo é dada pela equação:

\begin{equation}
  M_{\text{tubo}} = L_t \cdot m_{\text{tubo}}
  \label{eq:massa_tubo}
\end{equation}

Onde:

\begin{itemize}[leftmargin=1.25cm]
  \item $M_{\text{tubo}}$: massa do tubo em $\si{\kilo\gram}$.
  \item $L_t$: comprimento do tambor, convertido para $\si{\meter}$ \eqref{eq:comprimento_do_tambor}.
  \item $m_{\text{tubo}}$: massa do tubo por metro em $\si{\kilo\gram/\meter}$ (Tabela \ref{tab:dados_do_tambor_adotado}).
\end{itemize}

\vspace{1em}

Substituindo os valores na equação \eqref{eq:massa_tubo}:

\begin{align*}
  M_{\text{tubo}} & = 2112 \cdot 10^{-3} \cdot 194{,}64 \\
  M_{\text{tubo}} & = \boxed{\SI{411{,}08}{\kilo\gram}}
\end{align*}

A área do flange é dada pela equação:

\begin{equation}
  A_f = \frac{\pi \cdot {D_t}^2}{4}
  \label{eq:area_flange}
\end{equation}

Onde:

\begin{itemize}[leftmargin=1.25cm]
  \item $A_f$: área de cada flange em $\si{\meter^2}$.
  \item $D_t$: diâmetro externo do tubo, convertido para $\si{\meter}$ \ref{tab:dados_do_tambor_adotado}.
\end{itemize}

\vspace{1em}

Substituindo os valores na equação \eqref{eq:area_flange}:

\begin{align*}
  A_f & = \frac{\pi \cdot {\left(355{,}6 \cdot 10^{-3}\right)}^2}{4} \\
  A_f & = \boxed{\SI{99{,}32e-3}{\meter^2}}
\end{align*}

A massa do flange é dada pela equação:

\begin{equation}
  M_{\text{chapa}} = 3 \cdot A_f \cdot m_{\text{chapa}}
  \label{eq:massa_flange}
\end{equation}

Onde:

\begin{itemize}[leftmargin=1.25cm]
  \item $M_{\text{chapa}}$: massa do flange em $\si{\kilo\gram}$.
  \item $A_f$: área do flange em $\si{\meter^2}$.
  \item $m_{\text{chapa}}$: massa da chapa por metro quadrado em $\si{\kilo\gram/\meter^2}$.
\end{itemize}

\vspace{1em}

Com base na tabela de propriedades do aço ASTM A36 (Figura \ref{fig:chapas_astm_a36_sae_1045}), adotou-se uma chapa de \SI{12{,}7}{\milli\meter} de espessura, cuja massa por metro quadrado é \SI{99{,}7}{\kilo\gram/\meter^2}.

Substituindo os valores na equação \eqref{eq:massa_flange}:

\begin{align*}
  M_{\text{chapa}} & = 3 \cdot 99{,}32 \cdot 10^{-3} \cdot 99{,}7 \\
  M_{\text{chapa}} & = \boxed{\SI{29{,}71}{\kilo\gram}}
\end{align*}

O peso do tambor é dado pela equação:

\begin{equation}
  G_{\text{tambor}} = \left(M_{\text{tubo}} + M_{\text{chapa}}\right) \cdot g \cdot 1{,}5
  \label{eq:peso_tambor}
\end{equation}

Onde:

\begin{itemize}[leftmargin=1.25cm]
  \item $G_{\text{tambor}}$: peso do tambor em $\si{\newton}$.
  \item $M_{\text{tubo}}$: massa do tubo em $\si{\kilo\gram}$.
  \item $M_{\text{chapa}}$: massa total das chapas dos flanges em $\si{\kilo\gram}$.
  \item $g$: aceleração da gravidade em $\si{\meter\per\second\squared}$.
\end{itemize}

\vspace{1em}

Substituindo os valores na equação \eqref{eq:peso_tambor}:

\begin{align*}
  G_{\text{tambor}} & = \left(411{,}08 + 29{,}71\right) \cdot 10 \cdot 1{,}5 \\
  G_{\text{tambor}} & = \boxed{\SI{6611{,}85}{\newton}}
\end{align*}
